\chapter{Penurunan Diagonalisasi Operator Densitas dan Fungsi Logaritma Operator}
\label{app:diagonalisasi_operator}

Operator densitas $\rho$ merupakan operator Hermitian sehingga memenuhi
\begin{equation}
    \rho^\dagger=\rho.
\end{equation}
Berdasarkan Teorema Spektral, setiap operator Hermitian memiliki himpunan nilai eigen real beserta vektor-vektor eigen ortonormal yang memenuhi persamaan eigen
\begin{equation}
    \label{eq:persamaan-eigen-operator-densitas}
    \rho |e_j\rangle
    =
    \lambda_j |e_j\rangle,
\end{equation}
dengan $\lambda_j$ merupakan nilai eigen dan $|e_j\rangle$ merupakan vektor eigen ke-$j$. Ortonormalitas seluruh vektor eigen dinyatakan sebagai
\begin{equation}
    \langle e_i|e_j\rangle
    =
    \delta_{ij}.
\end{equation}
Seluruh vektor eigen kemudian disusun sebagai kolom-kolom suatu matriks
\begin{equation}
    U=
    \begin{bmatrix}
        |e_1\rangle &
        |e_2\rangle &
        \cdots &
        |e_n\rangle
    \end{bmatrix}.
\end{equation}
Karena seluruh vektor eigen saling ortonormal, maka berlaku
\begin{equation}
    U^\dagger U
    =
    I.
\end{equation}

Dengan mengalikan operator densitas terhadap matriks $U$, maka diperoleh
\begin{equation}
\begin{split}
    \rho U
    &=
    \rho
    \begin{bmatrix}
        |e_1\rangle &
        |e_2\rangle &
        \cdots &
        |e_n\rangle
    \end{bmatrix}
    \\
    &=
    \begin{bmatrix}
        \rho|e_1\rangle &
        \rho|e_2\rangle &
        \cdots &
        \rho|e_n\rangle
    \end{bmatrix}.
\end{split}
\end{equation}
Sesuai dengan persamaan \eqref{eq:persamaan-eigen-operator-densitas}, diperoleh
\begin{equation}
\begin{split}
    \rho U
    &=
    \begin{bmatrix}
        \lambda_1|e_1\rangle &
        \lambda_2|e_2\rangle &
        \cdots &
        \lambda_n|e_n\rangle
    \end{bmatrix}
    \\
    &=
    \begin{bmatrix}
        |e_1\rangle &
        |e_2\rangle &
        \cdots &
        |e_n\rangle
    \end{bmatrix}
    \begin{pmatrix}
        \lambda_1 &0&\cdots&0\\
        0&\lambda_2&\cdots&0\\
        \vdots&\vdots&\ddots&\vdots\\
        0&0&\cdots&\lambda_n
    \end{pmatrix}.
\end{split}
\end{equation}
Matriks diagonal tersebut dapat dinotasikan sebagai
\begin{equation}
    D=
    \operatorname{diag}
    (\lambda_1,\lambda_2,\ldots,\lambda_n),
\end{equation}
agar diperoleh hubungan
\begin{equation}
    \rho U
    =
    UD.
\end{equation}
Kemudian kedua ruas dikalikan dengan $U^\dagger$ sehingga diperoleh
\begin{equation}
    \rho UU^\dagger
    =
    UDU^\dagger.
\end{equation}
Karena $UU^\dagger=I$, maka diperoleh bentuk diagonalisasi operator
\begin{equation}
    \label{eq:diagonalisasi_rho}
    \rho
    =
    UDU^\dagger.
\end{equation}
Untuk memperoleh bentuk logaritma operator densitas, terlebih dahulu ditinjau deret Maclaurin dari fungsi logaritma
\begin{align}
    \log_2(1+x)
    &=
    \sum_{n=1}^{\infty}
    \frac{(-1)^{n-1}}{n\ln2}
    x^n,
    \\
    \log_2(x)
    &=
    \sum_{n=1}^{\infty}
    \frac{(-1)^{n-1}}{n\ln2}
    (x-1)^n.
\end{align}

Karena operator densitas merupakan operator Hermitian yang telah didiagonalisasi sebagaimana ditunjukkan pada persamaan~\eqref{eq:diagonalisasi_rho}, maka deret Maclaurin fungsi logaritma dapat diterapkan pada operator tersebut melalui \textit{functional calculus}. Maka dari itu, diperoleh representasi deret logaritma operator
\begin{equation}
\label{eq:deret_log_operator}
    \log_2(\rho)
    =
    \sum_{n=1}^{\infty}
    \frac{(-1)^{n-1}}
         {n\ln2}
    (\rho-I)^n.
\end{equation}
Karena operator densitas telah didiagonalisasi sebagaimana disajikan pada persamaan \eqref{eq:diagonalisasi_rho},maka
\begin{equation}
\begin{split}
(\rho-I)^n
&=
(UDU^\dagger-I)^n
\\
&=
\left(
U(D-I)U^\dagger
\right)^n.
\end{split}
\end{equation}
Dengan menggunakan sifat $U^\dagger U=I$, diperoleh
\begin{equation}
\begin{split}
(\rho-I)^n
&=
U(D-I)
\underbrace{U^\dagger U}_{I}
(D-I)
\cdots
\underbrace{U^\dagger U}_{I}
(D-I)
U^\dagger
\\
&=
U(D-I)^nU^\dagger.
\end{split}
\end{equation}
Karena
\begin{equation}
D=
\begin{pmatrix}
\lambda_1&0&\cdots&0\\
0&\lambda_2&\cdots&0\\
\vdots&\vdots&\ddots&\vdots\\
0&0&\cdots&\lambda_n
\end{pmatrix},
\end{equation}
maka
\begin{equation}
(D-I)^n=
\begin{pmatrix}
(\lambda_1-1)^n&0&\cdots&0\\
0&(\lambda_2-1)^n&\cdots&0\\
\vdots&\vdots&\ddots&\vdots\\
0&0&\cdots&(\lambda_n-1)^n
\end{pmatrix}.
\end{equation}
Substitusi hasil tersebut ke persamaan~\eqref{eq:deret_log_operator} sehingga menghasilkan
\begin{equation}
\begin{split}
\log_2(\rho)
&=
\sum_{n=1}^{\infty}
\frac{(-1)^{n-1}}
{n\ln2}
U(D-I)^nU^\dagger
\\
&=
U
\left[
\sum_{n=1}^{\infty}
\frac{(-1)^{n-1}}
{n\ln2}
(D-I)^n
\right]
U^\dagger.
\end{split}
\end{equation}
Karena matriks $U$ dan $U^\dagger$ tidak bergantung pada indeks penjumlahan $n$, kedua matriks tersebut dapat dikeluarkan dari operator penjumlahan. Sementara itu, matriks di dalam tanda kurung tetap berbentuk diagonal sehingga fungsi logaritma cukup diterapkan pada setiap elemen diagonalnya.
\begin{equation}
\label{eq:log-rho-diagonal}
\log_2(\rho)
=
U
\begin{pmatrix}
\log_2(\lambda_1)&0&\cdots&0\\
0&\log_2(\lambda_2)&\cdots&0\\
\vdots&\vdots&\ddots&\vdots\\
0&0&\cdots&\log_2(\lambda_n)
\end{pmatrix}
U^\dagger.
\end{equation}

Untuk memperoleh bentuk eksplisit dari hasil perkalian pada persamaan di atas, ditinjau terlebih dahulu suatu matriks diagonal umum
\begin{equation}
A=
\operatorname{diag}
(a_1,a_2,\ldots,a_n),
\end{equation}
dengan $a_j$ merupakan suatu konstanta. Untuk memperoleh bentuk eksplisit dari hasil perkalian pada persamaan~\eqref{eq:log-rho-diagonal}, ditinjau terlebih dahulu suatu matriks diagonal umum
\begin{equation}
\begin{split}
UA
&=
\begin{bmatrix}
|e_1\rangle&
|e_2\rangle&
\cdots&
|e_n\rangle
\end{bmatrix}
\begin{pmatrix}
a_1&0&\cdots&0\\
0&a_2&\cdots&0\\
\vdots&\vdots&\ddots&\vdots\\
0&0&\cdots&a_n
\end{pmatrix}
\\
&=
\begin{bmatrix}
a_1|e_1\rangle&
a_2|e_2\rangle&
\cdots&
a_n|e_n\rangle
\end{bmatrix}.
\end{split}
\end{equation}
Sementara itu, konjugat Hermitian dari matriks $U$ diberikan oleh
\begin{equation}
U^\dagger=
\begin{bmatrix}
\langle e_1|\\
\langle e_2|\\
\vdots\\
\langle e_n|
\end{bmatrix}.
\end{equation}
Sehingga bentuk penuh dari perkalian matriks diagonal tersebut adalah
\begin{equation}
\begin{split}
UAU^\dagger
&=
\begin{bmatrix}
a_1|e_1\rangle&
a_2|e_2\rangle&
\cdots&
a_n|e_n\rangle
\end{bmatrix}
\begin{bmatrix}
\langle e_1|\\
\langle e_2|\\
\vdots\\
\langle e_n|
\end{bmatrix}
\\
&=
a_1|e_1\rangle\langle e_1|
+
a_2|e_2\rangle\langle e_2|
+\cdots+
a_n|e_n\rangle\langle e_n|
\\
&=
\sum_{j=1}^{n}
a_j
|e_j\rangle
\langle e_j|.
\end{split}
\end{equation}

Dengan mengambil
\[
a_j=\log_2(\lambda_j),
\]
maka berdasarkan persamaan \eqref{eq:log-rho-diagonal} diperoleh
\begin{equation}
\log_2(\rho)
=
\sum_{j=1}^{n}
\log_2(\lambda_j)
|e_j\rangle
\langle e_j|.
\end{equation}
Karena operator densitas bersifat semidefinit positif, seluruh nilai eigennya memenuhi
\[
0\le\lambda_j\le1,
\]
dengan
\[
\sum_{j=1}^{n}\lambda_j=1.
\]
Apabila terdapat nilai eigen $\lambda_j=0$, fungsi $\log_2(\lambda_j)$ memang tidak terdefinisi pada titik tersebut. Akan tetapi, dalam Entropi Von Neumann yang muncul adalah hasil kali $\lambda_j\log_2(\lambda_j)$ sehingga nilainya didefinisikan melalui limit. Akan tetapi, pada Entropi Von Neumann suku yang muncul adalah hasil kali $\lambda_j\log_2(\lambda_j)$ sehingga nilainya didefinisikan melalui limit
\begin{equation}
\lim_{x\rightarrow0^{+}}
x\log_2x
=
0,
\end{equation}
sehingga keberadaan nilai eigen nol tidak menimbulkan singularitas pada perhitungan Entropi Von Neumann. 
